\documentclass[12pt,a4paper]{article}

% Packages
\usepackage[margin=25mm]{geometry}
\usepackage[T1]{fontenc}
\usepackage{lmodern}
\usepackage{xcolor}
\usepackage{hyperref}

\hypersetup{
    colorlinks=true,
    linkcolor=blue,
    urlcolor=blue
}

\title{\textbf{Building the Structure of a LaTeX Document}}
\author{LaTeX Learning Module}
\date{\today}

\begin{document}

\maketitle

\begin{abstract}
This example introduces the basic parts of a LaTeX document. It demonstrates
how to choose a document class, set common options, create a title, write an
abstract, and organize information with section headings.
\end{abstract}

\tableofcontents
\newpage

\section{Understanding the Document Class}

Every LaTeX document begins by selecting a document class. The class determines
the general structure and formatting rules used by the document.

\begin{itemize}
    \item \texttt{article} is useful for short documents, assignments, and reports
    that do not need chapters.
    \item \texttt{report} is suitable for longer work that can be divided into
    chapters.
    \item \texttt{book} is designed for book-length projects with front matter
    and chapters.
\end{itemize}

The class is selected with the \verb|\documentclass| command:

\begin{verbatim}
\documentclass[12pt,a4paper]{article}
\end{verbatim}

The options in square brackets set the main font size and paper format.
Common font-size choices include \texttt{10pt}, \texttt{11pt}, and
\texttt{12pt}. A document can also use paper options such as \texttt{a4paper}
or \texttt{letterpaper}.

\section{The Preamble and Document Body}

A LaTeX source file has two main parts. The preamble appears before
\verb|\begin{document}| and contains the document class, package imports,
and settings. The body begins after \verb|\begin{document}| and contains
the material that will appear in the finished document.

\begin{verbatim}
\documentclass[11pt,a4paper]{article}
\usepackage[margin=25mm]{geometry}

\begin{document}
Your document content goes here.
\end{document}
\end{verbatim}

Commands that load packages or define document-wide settings normally belong
in the preamble.

\section{Adding a Title and Author Details}

The commands \verb|\title|, \verb|\author|, and \verb|\date| store the
information that appears in the title area. The command \verb|\maketitle|
prints that information.

\begin{verbatim}
\title{A Short Research Report}
\author{Jordan Lee}
\date{November 2026}
\maketitle
\end{verbatim}

If the date field is left empty, use \verb|\date{}| to omit it. The appearance
of the title area depends partly on the selected document class.

\section{Writing an Abstract}

An abstract gives readers a brief overview of the document. In the
\texttt{article} class, it is usually created with the \texttt{abstract}
environment.

\begin{verbatim}
\begin{abstract}
This paper describes the purpose, approach, and main
points of the work in a concise form.
\end{abstract}
\end{verbatim}

Keep the abstract focused on the subject and the main information readers
should know before reading the full document.

\section{Organizing Content with Headings}

LaTeX provides commands for arranging material into a clear hierarchy.
For an article, common levels include:

\begin{itemize}
    \item \verb|\section{...}| for a main topic
    \item \verb|\subsection{...}| for a part of that topic
    \item \verb|\subsubsection{...}| for a more detailed point
\end{itemize}

For example:

\begin{verbatim}
\section{Methods}
\subsection{Materials}
\subsubsection{Equipment}
\end{verbatim}

The available heading levels depend on the document class. The
\texttt{report} and \texttt{book} classes also support chapters.

\section{A Complete Starter Example}

The following short source brings the main ideas together:

\begin{verbatim}
\documentclass[12pt,a4paper]{article}
\usepackage[margin=25mm]{geometry}

\title{My First Structured Document}
\author{Jordan Lee}
\date{\today}

\begin{document}
\maketitle

\begin{abstract}
A short summary of the document goes here.
\end{abstract}

\section{Introduction}
Introduce the subject.

\section{Main Discussion}
Explain the important points.

\subsection{A Supporting Detail}
Add a more focused explanation.

\section{Conclusion}
Summarize the main message.
\end{document}
\end{verbatim}

\section{Key Points}

A well-organized LaTeX file starts with a suitable class, places setup commands
in the preamble, and keeps the actual content inside the document body.
Titles, abstracts, and heading commands help readers navigate the finished
work.

\end{document}
